\section{Concept}\label{concept}

\subsection{Product type}

KompozeR is a \textbf{web application} (a browser-based Single Page Application)
backed by a set of cooperating \textbf{web services}. The end user interacts
with a graphical interface; the business logic is provided by several
independent backend services exposed through a single API Gateway.

The product extends a classic e-commerce scenario with a domain-specific
feature: a guided, visual \emph{configurator} (referred to as the CAD service)
that lets a customer compose a modular shelving unit out of individual
components. The catalog is organised in three mutually incompatible product
families (Tondo, Quadro, Qube), and the configurator enforces the domain
constraint that components of different families cannot be mixed in the same
configuration.

\subsection{Use case}

\begin{itemize}
  \item \textbf{What} the software does: it lets users browse a component
  catalog, interactively build a valid shelving configuration while seeing an
  up-to-date preview and price, save configurations to their profile, add them
  to a cart, and place orders. It keeps saved configurations consistent with
  the catalog, notifying users in real time when a component they depend on
  changes price or becomes unavailable. A contextual chatbot assists users
  during configuration, and administrators manage the catalog and inspect
  order trends.

  \item \textbf{Where} the users are: users are geographically distributed and
  reach the system over the public Internet from their own devices. Servers are
  centralised (a single logical deployment), while clients are remote.

  \item \textbf{When and how frequently}: interaction is on-demand and
  session-based. A configuration session is a relatively long, interactive
  activity; administrative catalog updates happen occasionally but must
  propagate promptly to every affected user.

  \item \textbf{How they interact / devices}: through a web browser (desktop or
  mobile) talking to the SPA. Communication is mostly request/response over
  HTTP; real-time channels (WebSocket via Socket.IO) are used for notifications
  and the chatbot.

  \item \textbf{Data to store}: yes. The system persists users and sessions,
  the component catalog, configurations and their bill of materials (BOM),
  carts, orders, notifications and subscriptions, and chat sessions. Each kind
  of data is owned and stored by the service responsible for it
  (database-per-service).

  \item \textbf{Roles}: the system distinguishes multiple roles ---
  \texttt{GUEST}, \texttt{BASE} (authenticated customer), and \texttt{ADMIN}
  (catalog and back-office operator) --- with different permissions.
\end{itemize}

\subsection{Why distribution is needed}

Distribution in KompozeR is a design consequence of the domain, not an end in
itself. The main motivations are:

\begin{itemize}
  \item \textbf{Resource sharing}: a single, authoritative catalog and its
  price/availability information must be shared among many concurrent users and
  among several application contexts (configurator, cart, notifications,
  chatbot, reporting).

  \item \textbf{Evolvability and separation of concerns}: the problem naturally
  decomposes into distinct bounded contexts (authentication, catalog,
  configuration, cart, orders, notifications, chatbot, reporting). Realising
  each as an independently deployable service allows them to evolve, scale, and
  be tested in isolation.

  \item \textbf{Fault tolerance}: the configuration workflow is the core of the
  product. Its persistence is backed by a replicated MongoDB replica set that
  survives the loss of a node through automatic primary election, and the
  collaborative session state is protected through checkpoint/recovery.

  \item \textbf{Scalability}: the load is dominated by many independent,
  read-heavy configuration sessions. A stateless service tier behind a gateway
  can be scaled horizontally to absorb an increasing number of users and
  requests.
\end{itemize}

Geographic distribution of servers and heavy shared computation are \emph{not}
primary drivers: the deployment is logically centralised, and the interesting
distributed-systems challenges are consistency of shared catalog data,
event-driven propagation of changes, and fault tolerance of the core workflow.
